%----------------------------------------------------------------------
\section{Plan Portfolios and Ambiguous-Parse Packing}
\label{sec:portfolios}\footnote{Section-level Toulmin. \emph{Move}: replace the single-plan output convention with a Pareto-portfolio output, and show how the ``waste'' lanes of LCM packing (Remark \ref{rem:coarseness}) are the natural home for multiple plans executed concurrently. \emph{[DAF] comparison}: [DAF]'s SPPF retains ambiguity as a first-class feature rather than forcing a deterministic choice; the plan portfolio is the scheduling-level analogue of the SPPF, retaining alternative derivations that are cost-indistinguishable under the stated hardware profile. \emph{Residue instance}: the gap between a plan's cost and the budget's slack is the residue that licenses additional packed plans in the remaining lanes.}
%----------------------------------------------------------------------

\subsection{Pareto frontier extraction}

\begin{definition}[Pareto frontier of a workload]
  \label{def:pareto}
  The \emph{Pareto frontier} $\mathrm{PF}(W, G, B)$ is
  defined by Proposition \ref{prop:frontier-char}: the
  set of plans admissible under $B$ that are not
  operationalist-dominated (Definition
  \ref{def:op-dominance}) by any other plan. The
  frontier is thus derived from the admissibility
  discriminator family (Definition
  \ref{def:admissibility-family}), not imposed as a
  multi-objective-optimization convention; the witness
  of dominance is a budget in the asymmetric
  difference of admissibility sets.
\end{definition}

\begin{proposition}[Pareto extraction from the chart]
  \label{prop:pareto}
  \footnote{Class \textsf{AH}. Warrant: a chart-item
  dominance check implements the partial order of
  Definition \ref{def:op-dominance} on the partial
  plans encoded as chart items. Qualifier: assumes
  eager pruning preserves the frontier; equal-cost
  items are both retained because neither dominates
  the other (Remark in
  Section \ref{def:admissibility-family} ff.).
  Finite-$\kappa$ valid.}
  The Pareto frontier $\mathrm{PF}(W, G, B)$ is the
  set of items in the chart cell $\chi[0, N]$ whose
  production is the workload's output nonterminal,
  whose cost vector satisfies $\bar{c} \leq B$
  (admissibility), and for which no other chart item
  with the same output nonterminal has a strictly
  smaller cost vector (non-dominance).
\end{proposition}

\begin{proof}
  The cost-bounded completer of
  Section \ref{sec:parsing} rejects chart items whose
  cost vector exceeds $B$, so every surviving item is
  admissible. Dominance pruning (applied eagerly on
  insertion) retains an item iff no existing item
  strictly dominates it in the sense of Definition
  \ref{def:op-dominance}. Items with equal cost
  vectors do not dominate each other, so are both
  retained; their distinguishability, if any, comes
  from the semantic attributes $(E, r, d)$ or the
  underlying atom sequence. By Proposition
  \ref{prop:frontier-char}, the surviving items at
  $\chi[0, N]$ for the output production are exactly
  $\mathrm{PF}(W, G, B)$.
\end{proof}

\subsection{Portfolio output}

\begin{definition}[Plan portfolio]
  \label{def:portfolio}
  A \emph{plan portfolio} is a triple
  $(\mathcal{P}, \mathcal{L}, \mathcal{E})$ where:
  \begin{itemize}[nosep]
    \item $\mathcal{P}$ is a set of plans drawn from
      (a subset of) the Pareto frontier;
    \item $\mathcal{L}$ is a lane partition of the
      packed register, assigning disjoint lane
      ranges to each plan in $\mathcal{P}$;
    \item $\mathcal{E}$ is the extraction rule
      (identity, XOR, majority vote) that maps the
      final register state to a single answer.
  \end{itemize}
\end{definition}

A portfolio with $|\mathcal{P}| = 1$ is a single
plan with no packing; a portfolio with $|\mathcal{P}|
> 1$ is a concurrent execution of multiple plans on
disjoint lane ranges.

\subsection{Packing constraint: shared atom sequences}

A key constraint limits which plans can share a
dispatch: in SIMT hardware, all lanes of a warp
execute the same instruction stream. Plans can
therefore pack only if their atom sequences
\emph{at the kernel level} are identical --- where
the lanes differ is in the \emph{data} they consume
and the \emph{masks} they apply.

\begin{definition}[Atom-sequence equivalence]
  \label{def:atom-equiv}
  Two plans $P_1$ and $P_2$ are
  \emph{atom-sequence equivalent} if their
  linearized atom sequences (in the order the
  planner would dispatch them) are identical. They
  may differ in encoding, tile size, masks, and
  data, but not in the sequence of opcodes.
\end{definition}

\begin{theorem}[Parse-packing soundness]
  \label{thm:parse-packing}
  \footnote{Class \textsf{AHO}$_{\mathrm{exec}}$.
  Warrant: mask-disjointness (as in
  Propositions \ref{prop:fused-correctness},
  \ref{prop:ts-commute}) plus the
  execution-operationalist axiom
  (Definition \ref{def:exec-operationalist})
  licenses concurrent execution of plans whose
  costs cannot be distinguished on the stated
  hardware. Qualifier: requires atom-sequence
  equivalence (Definition \ref{def:atom-equiv}).
  Finite-$\kappa$ valid. Novelty: the application of
  SIMT lane packing to alternative parse derivations
  is to our knowledge new; warp-level lane
  partitioning in standard GPU programming packs
  different data through the same kernel, not
  different plans for the same data.}
  Let $\mathcal{P}$ be a set of atom-sequence
  equivalent plans from the Pareto frontier, and
  $\mathcal{L}$ a lane partition assigning disjoint
  lane ranges to each. Then the concurrent
  execution of all plans in $\mathcal{P}$ on the
  partitioned register produces each plan's output
  on its assigned lane range, and any one of those
  outputs is a valid result for the workload.
\end{theorem}

\begin{proof}
  By atom-sequence equivalence, all plans execute
  the same instruction stream. By lane disjointness,
  no plan's mask overlaps another's, so the
  block-diagonal argument of Proposition
  \ref{prop:fused-correctness} applies: each plan's
  output on its lane range depends only on its own
  data. By soundness (Theorem \ref{thm:soundness}),
  each plan's output is a correct result for the
  workload. Finally, since all plans satisfy the
  budget, the execution-operationalist axiom
  (Definition \ref{def:exec-operationalist}) makes
  any of them an acceptable answer.
\end{proof}

\subsection{The three extraction modes}

\begin{definition}[Extraction modes]\label{def:extract}
  A portfolio's extraction rule $\mathcal{E}$ is
  one of:
  \begin{description}[nosep,leftmargin=1cm]
    \item[\textsc{First}] Read the output of any
      single plan from its assigned lane range. The
      simplest mode; used when all plans are
      cost-indistinguishable.
    \item[\textsc{XOR}] For an odd number of plans,
      XOR their outputs: over $\GF(2)$, the XOR of
      $k$ agreeing values is the value itself when
      $k$ is odd. Detects single-bit disagreements
      (useful for fault detection when alternative
      plans are cost-distinguishable but all
      supposed to compute the same value).
    \item[\textsc{Majority}] For 2 or more plans,
      take the majority vote lane-by-lane.
      Triple-modular-redundant (TMR) in the sense
      used by reliability engineering. Tolerates
      single-plan failure.
  \end{description}
\end{definition}

\begin{remark}[Batch coarseness as a speculation budget]
  Proposition \ref{prop:lcm-density} showed that
  under LCM packing the hardware tile's
  information density is constant at $M^2$ bits,
  which forces a minimum batch size $(L/d)^2$. If
  the workload's batch is smaller than this minimum,
  the remaining lanes would otherwise be wasted
  padding; portfolio packing uses them for
  alternative plans at zero marginal cost.
\end{remark}

\subsection{Enumeration strategies}

Three strategies for choosing which plans to pack:

\begin{itemize}[nosep]
  \item \textsc{Top-$k$}: pack the $k$ lowest-cost
    plans from the Pareto frontier, weighted by
    some scalar $\phi$. Maximizes the probability
    that the realized-best plan is among the
    packed.
  \item \textsc{Corner-pick}: pack one plan that
    minimizes each of the $d$ cost dimensions
    separately (so $d$ plans total). Ensures that
    whichever resource turns out to be binding in
    practice has a dedicated plan optimized for it.
  \item \textsc{Ambiguous-cluster}: group
    atom-sequence-equivalent plans and pack the
    entire cluster. Useful when the grammar admits
    many plans that differ only in encoding or
    tile assignment.
\end{itemize}

Each strategy is a wrapper around the Pareto
frontier extraction of Proposition \ref{prop:pareto};
all three produce valid portfolios under
Theorem \ref{thm:parse-packing}.

\subsection{Plan-equivalence transitivity}

\begin{proposition}[Plan-equivalence transitivity]
  \label{prop:plan-trans}
  \footnote{Class \textsf{AO}$_{\mathrm{exec}}$.
  Warrant: direct from
  Definition \ref{def:exec-operationalist}.
  Finite-$\kappa$ valid.}
  Let $P_1, P_2, P_3$ be plans for the same
  workload on the same hardware. If $P_1 \equiv
  P_2$ and $P_2 \equiv P_3$ (execution-equivalent
  per Definition \ref{def:exec-operationalist}),
  then $P_1 \equiv P_3$.
\end{proposition}

\begin{proof}
  Measurement indistinguishability is transitive:
  if no observation separates $P_1$ from $P_2$ and
  no observation separates $P_2$ from $P_3$, no
  observation separates $P_1$ from $P_3$.
\end{proof}

This makes $\equiv$ a genuine equivalence relation
on plans, and licenses the parser's portfolio
output to be a single representative per equivalence
class rather than an enumeration of duplicates.

\subsection{Register-time-space allocation}

\begin{theorem}[Register-time-space bound]
  \label{thm:reg-bound}
  \footnote{Class \textsf{AH}. Warrant: interval-graph
  coloring argument on the lane $\times$ time
  allocation. Qualifier: assumes register lanes and
  dispatch cycles form a 2D grid with rectangular
  claims. Finite-$\kappa$ valid.}
  For a portfolio with $k$ plans, each requiring
  register width $w_i$ and dispatch-cycle count
  $t_i$, a valid packing requires register-time
  area at least $\sum_i w_i \cdot t_i$; the greedy
  strip-packing allocation attains area at most
  $2 \sum_i w_i \cdot t_i$.
\end{theorem}

\begin{proof}
  The lower bound is a straightforward counting
  argument: total area $\geq$ sum of per-plan
  areas. The greedy upper bound is the standard
  strip-packing result (Coffman, Garey, Johnson,
  Tarjan 1980), which is 2-approximate for
  rectangular strip packing.
\end{proof}

\subsection{Summary of the planner's outputs}

The planner, in its most general form, takes a
workload $W$, a grammar $G$, a budget $B$, a
hardware profile $\pi$, and an objective $\phi$,
and returns:

\begin{itemize}[nosep]
  \item A plan portfolio
    $(\mathcal{P}, \mathcal{L}, \mathcal{E})$
    satisfying Theorem \ref{thm:parse-packing};
  \item The dual prices $(\lambda_{\VRAM}^*,
    \lambda_{\SMEM}^*, \lambda_{\REG}^*)$ from
    the Lagrangian dual
    (Section \ref{sec:lagrangian}), interpretable
    as marginal per-resource shadow costs;
  \item A diagnostic report on which resource
    is binding, alternative plans within a stated
    cost tolerance, and the estimated Pareto
    frontier size.
\end{itemize}

This is a substantially richer output than the
single-plan output of a traditional compiler. It
enables the end user to reason about the tradeoff
surface directly rather than about a pre-collapsed
scalar answer. Sections \S1 S1--S2 of the reading
guide note this as the central engineering
investment required to make the planner
operationally useful; the grammar and parser
described in this paper are the infrastructure
that makes such an investment pay off.
